\begin{corollary}[Two far inputs and almost-certain nearby mixtures]
\label{pd:two-far-inputs}
In Corollary~\ref{pd:full-gap}, let $r\ge2$. When $r$ is even,
there are words $u,v$ with
\[
 \dist(u,C)=\dist(v,C)=\theta+\frac{r}{2r+1}\eta
\]
such that
\[
 \dist((1-z)u+zv,C)=\theta
 \qquad\text{for every }z\in\F_p\setminus\{0,1\}.
\]
Thus exactly the two endpoints are outside the tested radius, and a
uniform affine mixture is nearby with probability $1-2/p$.
For $r\to\infty$, both endpoints are almost half a gap outside that
radius, although the numerical prescription predicts a nearby fraction
tending to zero. All the same rate, domain-length, strict-Elias, and
sampling guarantees apply.

More generally, for every divisor $t\ge2$ of $r$, the same code has
$t$ words $v_1,\ldots,v_t$, each at distance
\[
 \theta+\frac{2r}{2r+1}\left(1-\frac1t\right)\eta,
\]
whose uniform affine combination is nearby with exact probability
\[
 \left(1-\frac1p\right)^t-\frac{(-1)^t}{p^t}
 =1-O(t/p).
\]
Here the coefficients are uniform subject to summing to one.
For each fixed batch size $t$, taking $r\to\infty$ through multiples
of $t$ makes every input almost $(1-1/t)\eta$ outside the radius.
\end{corollary}
\begin{proof}
For the first statement, partition the $r$ directions into two equal
sets $A,B$ and put
$u=f+\sum_{j\in A}e_j$, $v=f+\sum_{j\in B}e_j$.
Their affine mixture has parameter coordinates $1-z$ on $A$ and
$z$ on $B$. At the endpoints exactly $r/2$ coordinates are nonzero;
at every other parameter all $r$ are nonzero.
The exact distance profile~\eqref{pd:exact-profile} proves every claim
about the line. Its $p-2$ nearby parameters still exceed the prescribed
bound by a factor tending to infinity.

For the second statement partition the directions into $t$ groups
$A_j$ of size $r/t$ and take $v_j=f+\sum_{i\in A_j}e_i$.
Each word has parameter weight $r/t$, giving the stated distance.
An affine combination has coefficient $\lambda_j$ repeated on $A_j$,
so it is nearby exactly when all $t$ coefficients are nonzero.
The number of nonzero-coordinate vectors with sum one is
$((p-1)^t-(-1)^t)/p$: this follows by inclusion--exclusion, or by
separating zero and nonzero target sums in a recurrence on $t$.
Dividing by the $p^{t-1}$ affine coefficient vectors proves the formula.
\end{proof}

\begin{corollary}[Linear and multilinear mixtures of far inputs]
\label{pd:far-batching}
Use the $t$ words of Corollary~\ref{pd:two-far-inputs} and assume
$n-K\ge4r+1$, as holds throughout its sufficiently large-field
asymptotic regime. A uniform unrestricted linear combination is
nearby with exact probability
\[
 \frac1p+\left(1-\frac1p\right)
 \left[\left(1-\frac1p\right)^t-\frac{(-1)^t}{p^t}\right].
\]
If $t=2^\ell$, index the words by $b\in\{0,1\}^\ell$ and use the
multilinear coefficients
\[
 \lambda_b(x)=\prod_{j:b_j=0}(1-x_j)\prod_{j:b_j=1}x_j.
\]
For uniform $x\in\F_p^\ell$, the mixture is nearby with exact
probability $(1-2/p)^\ell$. For every fixed $\ell$, all inputs can
lie almost $(1-2^{-\ell})\eta$ outside the tested radius while this
probability tends to one.
\end{corollary}
\begin{proof}
Let $s$ be the sum of the linear coefficients. When $s\ne0$,
scaling by $s^{-1}$ preserves distance to the linear code and produces
a uniform affine combination, with the probability already computed.
When $s=0$, the $f$ terms cancel and the mixture is supported on at
most $2r$ coordinates. Its distance to the zero codeword is at most
$2r/n\le\theta$. Since $s$ is uniform, the first formula follows.
For the multilinear statement the coefficients sum to one. All are
nonzero exactly when every $x_j$ avoids zero and one. Apply the exact
affine profile and independence of the $\ell$ parameters.
\end{proof}
